\subsection{Classical extraction from the concrete transcript}\label{app:concrete-extraction}
This subsection analyzes the Q version 2 and G version 1 profiles in a classical random-oracle model. All oracle stages refer to Appendix~\ref{app:namespace-composition}. Counts include adversarial queries, honest and simulated services, all four scalar lanes of a vectorized call, and final verification. The extractor sees the controller's private query log; that log is not part of the adversary's public view. On first touching an XOF input, the controller caches the full protocol-length prefix (768 bits at stage 6 and the required lengths elsewhere), revealing only the prefix requested by the caller.

Both profiles use $\tau=22$ small-VOLE vectors: eight have $4096$ leaves and fourteen have $2048$, totaling $L=61{,}440$. Their hidden indices consume 250 challenge bits; the remaining six bits must be zero. Stage 6 maps a leaf seed and IV to a 256-bit message and a 512-bit commitment. Stage 7 hashes each ordered leaf-commitment vector to 512 bits; stage 8 hashes the 22 vector digests to a 512-bit commitment digest. Stage 9 hashes the 256 ordered 272-bit column checks to a 512-bit digest.

\paragraph{Commitment extraction.}
At the first relevant stage-3 query, snapshot the logs. Follow the commitment digest through its first correctly sized stage-8 preimage, each stage-7 preimage, and each stage-6 commitment suffix with the prescribed IV. Output the corresponding leaf messages, retaining $\bot$ wherever a preimage is absent. This is a table lookup, not inversion of an unqueried hash.

The compressed opening reveals seeds at the frontier of the union of the 22 hidden paths. The verifier expands every unhidden leaf and recomputes its stage-6 commitment, then stages 7 and 8. Thus an accepted opening can disagree with an extracted message only through a digest collision, or by later hitting a digest that had no preimage at the snapshot. Compression changes the encoding of the opened seeds, not this argument. If $S_3$ snapshots and $q_a$ distinct stage-$a$ inputs occur, a conservative bound is
\[
 \eps_{\rm VC}\le
 \frac{\sum_{a=6}^8\binom{q_a}{2}
       +S_3(q_8+22q_7+61{,}440q_6)}{2^{512}}.
\]
The bound permits repeated leaf commitments and charges ambiguities to collisions. The threshold on opening length only rejects additional openings. A stage-9 log lookup at each stage-4 snapshot similarly extracts all column-check values with failure at most
\[
 \eps_9\le[\tbinom{q_9}{2}+S_4q_9]/2^{512}.
\]

\paragraph{The implemented VOLE check.}
Let $\ell$ be the number of QuickSilver rows, and retain exactly $\ell+272$ bits of each expanded stream; allocation padding is excluded. Put $n=\lceil\ell/256\rceil$, which is 56 for G and 58 for Q. The actual check is a binary linear map $\mathcal L_\chi$ consisting of a full-field polynomial hash, a 64-bit polynomial hash, a random two-by-two field matrix, truncation to 272 bits, and XOR with the 272-bit tail mask. For a fixed nonzero difference $e$ and fresh $\chi$,
\[
 \Pr[\mathcal L_\chi(e)=0]\le (n-1)/2^{256}+2^{-272}.
\]
Indeed, the full-field data polynomial has at most $n-1$ roots. Otherwise the two independent matrix rows supply a uniform 256-bit first output and 16 retained bits from the second. A difference confined to the tail mask never vanishes.

\paragraph{Extracting one affine witness.}
For a completely extracted vector $i$, let $z_{i,x}$ be its expanded leaf streams, $c_0=0$, and $c_i$ its transmitted correction. Define
\[
 u_i=\bigoplus_x z_{i,x}\oplus c_i,\qquad
 V_{i,t}=\bigoplus_{x:x_t=1}z_{i,x}.
\]
The receiver's XOR and Gray-code conversion gives
$Q_{i,t}=V_{i,t}\oplus\delta_{i,t}u_i$.
Before stage 4, extract the declared column values $D_{i,t}$ from stage 9 and read the transmitted check $\widehat u$. An accepting continuation must satisfy
\[
 D_{i,t}=\mathcal L_\chi(V_{i,t})\oplus
 \delta_{i,t}\bigl(\mathcal L_\chi(u_i)\oplus\widehat u\bigr).
\]
Call a complete vector an anchor when $\mathcal L_\chi(u_i)=\widehat u$. The $u_i$ are fixed before $\chi$; except with probability $\binom{22}{2}((n-1)2^{-256}+2^{-272})$ per fresh challenge, equal check values among complete vectors imply equal relevant row strings.

If an anchor exists, choose its $u$. Every complete nonanchor block has its entire $\delta_i$ fixed by the displayed equation, or cannot accept. For its required value $\beta_i$, replace its extractor matrix by $V^*_{i,t}=V_{i,t}\oplus\beta_{i,t}(u_i\oplus u)$. An incomplete vector with two missing leaves cannot open; with exactly one missing leaf $h$, acceptance fixes $\delta_i=h$. Fill that missing \emph{row string} with $u\oplus c_i\oplus\bigoplus_{x\ne h}z_{i,x}$ for the extractor's algebra only. Its coefficient in the receiver computation is zero, so no PRG preimage is needed. These operations produce one $V^*,u$, fixed before stage 4, such that every accepting continuation satisfies $Q=V^*\oplus u\Delta$. The witness correction $d$ then gives the Boolean witness $w=d\oplus u_W$.

If no anchor exists, every complete block fixes its challenge bits and every incomplete block fixes its hidden index. Including the six required zeros, at most one 256-bit final challenge can accept. This costs at most $q_5/2^{256}$ over raw fresh final-stage queries. It is not a probability conditioned on successful grinding.

\paragraph{Fresh chains after simulation.}
The public simulator in Appendix~\ref{app:public-simulation} programs only stage 5 on its own complete statement chains. Stages 1, 3, 4 and 6--9 remain unprogrammed. For a fixed kind and profile, the stage-1 encoding of the complete public statement and context is injective; collisions in application-level context digests are accounted for separately. Their chained outputs have lengths $b_1=512$, $b_3=1344$, and $b_4=832$ for G or $1664$ for Q. Collisions and a dependent input guessed before its earlier preimage is queried have combined probability at most
\begin{align*}
 \eps_{\rm chain}\le{}&
 \frac{\binom{q_1}{2}+q_1q_3}{2^{512}}+
 \frac{\binom{q_3}{2}+q_3q_4}{2^{1344}}+
 \frac{\binom{q_4}{2}+q_4q_5}{2^{b_4}}.
\end{align*}
Outside this event, a final proof determines a unique chronological chain. Sharing a programmed stage-5 input forces the same stage-4 output and preimage, then the same stage-3 preimage, and finally the same statement digest and public statement. A fresh statement therefore cannot use a simulated programming point.

\begin{theorem}[Classical single-output extraction]\label{thm:concrete-wse}
For the stated profile and the matched controller of Appendix~\ref{app:public-simulation}, there is a straight-line extractor for a fresh accepted full statement with errors bounded by
\begin{align*}
 e_Q\le{}&\eps_{\rm VC}+\eps_9+\eps_{\rm chain}
 +231q_3(57\cdot2^{-256}+2^{-272})\\
 &+q_4(M+F)2^{-256}+5q_5\,2^{-256},\\
 e_G\le{}&\eps_{\rm VC}+\eps_9+\eps_{\rm chain}
 +231q_3(55\cdot2^{-256}+2^{-272})\\
 &+3688q_4\,2^{-256}+3q_5\,2^{-256}.
\end{align*}
Each bound is capped at one. Parameter-digest and outer application collisions are accounted for separately. The quantities $M,F$ are those of Lemma~\ref{lem:batching}.
\end{theorem}
\begin{proof}
Verify the final object and follow its logged chain. Use the pre-stage-3 snapshot for leaf extraction and the pre-stage-4 snapshot for the check values and affine normalization. If an anchor exists, test $w=d\oplus u_W$ against the encoded relation and output its key, and its selector index for Q. This is polynomial-time in the query count, configured vector size and row length.

Outside the digest and ordering events, an invalid extracted witness survives only a VOLE-check collision, a residual batching erasure, or a root of the final polynomial. The first costs the stated pairwise term. Lemma~\ref{lem:batching} bounds the second. The final polynomial has degree four for Q or two for G, contributing $4q_5/2^{256}$ or $2q_5/2^{256}$. Adding the no-anchor event gives the displayed final coefficients. A fresh statement uses an unprogrammed final point by the chain argument, so the same bounds hold after adaptive simulation queries. Only the final selected witness is returned.
\end{proof}

